\section{What the Model Teaches}

The Kuramoto model is deliberately small: it keeps only phase, natural frequency, and coupling. Yet it
explains several general lessons.
\begin{itemize}[label=$\diamond$]
    \item A shared rhythm can emerge without a leader or global clock.
    \item The same mathematical language can describe biological rhythms, power-grid generators, chemical oscillations, and distributed timing systems.
\end{itemize}

The model also has limitations. Real fireflies interact through visible pulses, have restricted fields of
view, move through space, and may respond differently to males and females. Those details require
richer pulse-coupled or spatial models. A field study and video of natural synchronous swarms are
available through the \href{https://doi.org/10.1126/sciadv.abg9259}{Science Advances study by Sarfati, Hayes, and Peleg.}